Báo cáo này tập trung vào việc nghiên cứu và triển khai hai thuật toán trí tuệ nhân tạo (AI) phổ biến trong trò chơi cờ vua: Alpha-Beta Pruning và Monte Carlo Tree Search (MCTS). Mục tiêu chính là xây dựng một hệ thống AI có khả năng đưa ra các quyết định tối ưu dựa trên các phương pháp tìm kiếm và đánh giá khác nhau.

Hệ thống được thiết kế để so sánh hiệu quả giữa phương pháp duyệt cây tìm kiếm có hệ thống và phương pháp mô phỏng xác suất. Alpha-Beta Pruning, một cải tiến của thuật toán Minimax, giúp giảm thiểu đáng kể không gian tìm kiếm bằng cách loại bỏ các nhánh không khả thi thông qua hai chỉ số Alpha và Beta. Để tối ưu hóa quyết định, thuật toán sử dụng hàm đánh giá Heuristic kết hợp giữa giá trị quân cờ (Material Score) và điểm vị trí (Position Square Table).

Song song đó, thuật toán Monte Carlo Tree Search (MCTS) tiếp cận bài toán thông qua việc mô phỏng hàng ngàn ván đấu ngẫu nhiên để đưa ra quyết định dựa trên thống kê. Quy trình của MCTS bao gồm bốn giai đoạn cốt lõi: Chọn lọc (Selection) dựa trên công thức UCB1 để cân bằng giữa khai thác và khám phá, Mở rộng (Expansion), Mô phỏng (Simulation) và Truyền ngược (Backpropagation). 

Kết quả nghiên cứu cung cấp cái nhìn chi tiết về ưu và nhược điểm của từng phương pháp: Alpha-Beta mạnh mẽ trong việc tìm nước đi tối ưu ở độ sâu nhất định, trong khi MCTS linh hoạt và hiệu quả trong các không gian trạng thái khổng lồ mà không phụ thuộc quá nhiều vào hàm đánh giá thủ công. Đây là nền tảng quan trọng để phát triển các hệ thống AI chơi cờ chuyên sâu hơn.
